\documentclass{article}
\usepackage{amsmath, amssymb, amsfonts}
\usepackage{hyperref}
\usepackage{graphicx}
\usepackage{geometry}
\geometry{a4paper, margin=1in}

\title{Mathematical Derivations and Code Mappings for BrainKAN Explainability}
\author{Xin Qi}
\date{October 2026}

\begin{document}
\maketitle

\begin{abstract}
This document outlines the formal mathematical framework for the BrainKAN explainability project, providing one-to-one mappings between the theoretical derivations and their concrete implementations in the Python codebase. Crucially, it formulates rigorous statistical null-models (permutation testing) to demonstrate that many apparent non-linear computational effects in macro-scale connectomics are artifacts of input distribution shifts, highlighting the necessity of structural null-calibration.
\end{abstract}

\section{Edge-Specific KAN Convolution Architecture}
\subsection{Mathematical Formulation}
Let the task-fMRI functional brain network be represented as a directed graph $G = (V, E)$, where $V$ represents the brain regions of interest (ROIs) and $E$ represents the edges (functional connectivities). For a given layer $l$, the node features are updated via:
\begin{equation}
h_j^{(l+1)} = \sum_{i \in \mathcal{N}(j)} \Phi_{ij}^{(l)}\left(h_i^{(l)}\right)
\end{equation}
Unlike conventional message-passing layers with shared transformations, BrainKAN parameterizes $\Phi_{ij}$ uniquely for every existing edge $(i,j) \in E$ using B-splines. 

For a scalar input $x$, the edge-specific function is defined as:
\begin{equation}
\Phi_{ij}(x) = W^{base}_{ij}\operatorname{SiLU}(x) + \operatorname{spline}_{ij}(x) = W^{base}_{ij} \operatorname{SiLU}(x) + \sum_{k=1}^K w_{ij,k} B_k(x)
\end{equation}
where $B_k(x)$ are the B-spline basis functions defined recursively over a uniform grid, and $w_{ij,k}$ are the learnable spline coefficients specifically allocated for the edge $(i, j)$.

The updated node features are:
\begin{equation}
h_i' = R(h_i) + \sum_{j} A_{ij}\Phi_{ij}(h_j)
\end{equation}
where $R$ is an optional residual connection.

\subsection{Code Implementation}
\textbf{File:} \texttt{models/kan.py} \\
\textbf{Class/Function:} \texttt{EdgeSpecificKANConv.forward()} and \texttt{\_b\_spline\_basis()} \\
The recursive B-spline basis evaluation is centrally implemented in \texttt{\_b\_spline\_basis()}, which is imported throughout the repository to guarantee numerical consistency (e.g., in \texttt{02\_q3b\_common\_reference\_geometry.py}).

\section{Q1: Function Characterization and Deviation-from-Linearity (DFL)}
\subsection{Mathematical Formulation}
To characterize whether the learned edge function $\Phi_{ij}$ is actually nonlinear, we evaluate it over a continuous domain $T = [t_{min}, t_{max}]$ and compute the Deviation-From-Linearity (DFL). Let $y(t) = \Phi_{ij}(t)$. We fit a simple linear regression $\hat{y}(t) = a t + b$.
The DFL is defined as $1 - R^2$, where:
\begin{equation}
R^2 = 1 - \frac{\sum_{t \in T} (y(t) - \hat{y}(t))^2}{\sum_{t \in T} (y(t) - \bar{y})^2}
\end{equation}
\begin{equation}
\operatorname{DFL} = 1 - R^2 = \frac{\sum_{t \in T} (y(t) - \hat{y}(t))^2}{\sum_{t \in T} (y(t) - \bar{y})^2}
\end{equation}
A higher DFL indicates that the edge function deviates significantly from a pure linear transformation.

\subsection{Code Implementation}
\textbf{File:} \texttt{01\_q1\_q2\_q3a\_naive\_analysis.py} \\
\textbf{Function:} \texttt{compute\_function\_signatures()} \\
The code computes \texttt{ss\_res} (residual sum of squares) and \texttt{ss\_tot} (total sum of squares) to derive $R^2$, and then outputs $\operatorname{DFL}_e = 1 - R^2$. The edge-level DFL is averaged across output channels: $NL_e = \frac{1}{16}\sum_o NL_{e,o}$.

\section{Q2: Function Taxonomy}
\subsection{Mathematical Formulation}
To determine if edge functions group into a discrete taxonomy of biological computations, we standardize the function signatures and apply Ward hierarchical clustering for $k=2, \dots, 7$. Cluster quality and fold-level cross-validation stability are measured using Silhouette and Adjusted Rand Index (ARI):
\begin{equation}
\operatorname{Silhouette} \approx 0.18 - 0.19, \quad \operatorname{ARI} \approx 0
\end{equation}
\textit{This indicates no reproducible small-$k$ discrete taxonomy is detected across folds.}

\subsection{Code Implementation}
\textbf{Files:} \texttt{01\_q1\_q2\_q3a\_naive\_analysis.py} (Real) and \texttt{05\_q2\_null\_taxonomy.py} (Null).

\section{Q3A: Naive Curvature Remodeling}
\subsection{Mathematical Formulation}
The naive curvature effect ($\Delta C$) assesses whether the second derivative of the edge function changes across task conditions ($c \in \{0BK, 2BK\}$). For an edge $e$, the condition-specific naive curvature is calculated over the subject-specific input distribution $X^{(c)}$:
\begin{equation}
C_e(x) = \frac{1}{16} \sum_{o=1}^{16} |f_{e,o}''(x)|
\end{equation}
The shift in curvature is then:
\begin{equation}
\Delta C_e = C_e(x_j^{2BK}) - C_e(x_j^{0BK})
\end{equation}

\textit{Note: As discussed, $\Delta C_e$ is mathematically confounded by kinematic shifts. Locally, away from points where $f''(x)=0$, a first-order expansion illustrates how operating-regime shifts can induce apparent changes in curvature statistics:}
\begin{equation}
\Delta C \approx C'(x_0) \cdot \Delta x + \frac{1}{2}C''(x_0) \cdot \Delta x^2 + \cdots
\end{equation}
\textit{This demonstrates why $\Delta C^{real} \approx \Delta C^{null}$ under a label-shuffled null, representing an identifiability failure of the naive estimand.}

\subsection{Code Implementation}
\textbf{File:} \texttt{01\_q1\_q2\_q3a\_naive\_analysis.py} \\
The script evaluates the second derivative using PyTorch's \texttt{torch.autograd.grad} on the actual task-specific inputs.

\section{Q3B: Common-Reference Null Calibration}
\subsection{Mathematical Formulation}
To eliminate the input-regime confounder, we evaluate the learned geometry over a fixed, condition-independent reference grid $x^{ref}_j$, constructed via quantiles of the pooled training distribution:
\begin{equation}
x^{ref}_j = Q_{\text{pooled training},j}(q)
\end{equation}
The intrinsic structural nonlinearity (using DFL as the test statistic $G_e$) is computed for the real model ($G^{real}$) and compared against $R=100$ null models ($G^{null}$):
\begin{equation}
G_e = \frac{1}{16} \sum_o (1 - R^2_{e,o,\text{ref}})
\end{equation}
The observed statistic is:
\begin{equation}
T_e^{real} = G_e^{real}
\end{equation}
and compared against $\{G_e^{null,r}\}_{r=1}^R$. 
The effect size measuring excess geometry is:
\begin{equation}
\Psi_e = G_e^{real} - E[G_e^{null}]
\end{equation}
$\Psi_e$ represents the true biological signal isolated from structural and distributional artifacts. Statistical significance (yielding 0/240 FDR) is determined by whether $T_e^{real}$ lies in the extreme tail of the empirical null distribution.

\subsection{Code Implementation}
\textbf{File:} \texttt{02\_q3b\_common\_reference\_geometry.py} \\
Models are trained identically under true labels vs. subject-level swapped labels. All functions are probed strictly on $X_{ref} = [-8, 12]$.

\section{Q3C: Held-Out Calibrated Effective Gain}
\subsection{Mathematical Formulation}
Because Q3B found no excess structural nonlinearity, we test for simpler computational changes: changes in the Effective Gain (the first derivative). To ensure unbiased estimation, we use a fixed condition-independent reference grid $X_{ref}$, strictly computed from the \textit{training fold} to prevent data leakage.
\begin{equation}
G_{e,c} = \frac{1}{16} \sum_o \mathbb{E}_{x \sim P_c} \left[ |f_{e,o}'(x)| \right]
\end{equation}
Since $f'(x)$ is evaluated via numerical differentiation:
\begin{equation}
f'(x) \approx \frac{f(x+\epsilon) - f(x-\epsilon)}{2\epsilon}
\end{equation}
The paired difference in effective gain for a subject is:
\begin{equation}
dG_i = G_{i}^{2BK} - G_{i}^{0BK}
\end{equation}
We test the hypothesis that the paired differences from the real model differ from the null:
\begin{equation}
H_0 \Rightarrow (dG_i^{real}, dG_i^{null}) \overset{d}{=} (dG_i^{null}, dG_i^{real})
\end{equation}
Let $T^{real}_e = \sum_i (dG_i^{real} - dG_i^{null})$ and for each permutation $r$, the null statistic is $T^{null,r}_e = \sum_i s_i^{(r)} (dG_i^{real} - dG_i^{null})$ with $s_i^{(r)} \in \{-1, +1\}$. The permutation p-value with $R$ sign-flips is:
\begin{equation}
p_e = \frac{1 + \sum_r I(|T^{null,r}_e| \ge |T^{real}_e|)}{R + 1}
\end{equation}
\begin{equation}
d_z = \frac{\mu_{\Delta G}}{\sigma_{\Delta G}}
\end{equation}
where $\mu_{\Delta G}$ and $\sigma_{\Delta G}$ are the mean and standard deviation of $dG_{e,s}^{real} - dG_{e,s}^{null}$ across all test subjects.

\subsection{Code Implementation}
\textbf{File:} \texttt{04\_q3c\_effective\_gain.py} \\
\textbf{Functions:} \texttt{evaluate\_Q3C()} and the subsequent statistical testing block. \\
Cohen's $d_z$ is explicitly calculated as \texttt{observed\_mean / np.std(diff)} and saved in \texttt{q3c\_final\_table.csv}.

\section{Synthetic Identifiability Benchmark}
\subsection{Mathematical Formulation}
To define the identifiability boundary, we compare the measured effective gain $\Delta G$ across synthetic scenarios: baseline ($S_0$), input shift ($S_1$), and gain modulation ($S_2$).
\begin{equation}
S_1 (\text{input shift}): \Delta G = -0.087 \pm 0.13
\end{equation}
\begin{equation}
S_2 (\text{gain modulation}): \Delta G = -0.082 \pm 0.13
\end{equation}
\begin{equation}
S_1 \approx S_2
\end{equation}
This demonstrates empirically that input shift and gain modulation are confusable under the tested Q3C estimator. To ensure the estimator can detect true structural remodeling when present, $S_5$ tests a strong true structural change via a common-support comparison distance $D$:
\begin{equation}
D_{S5} = 3.574 \pm 0.75 \gg D_{S0} = 0.153 \pm 0.25
\end{equation}
Rejecting $H_0: S_5 = S_0$ proves the estimator responds to the tested remodeling regime.

\subsection{Code Implementation}
\textbf{File:} \texttt{03\_synthetic\_identifiability\_benchmark.py} \\
Simulates S0-S5 using synthetic 1D splines.

\end{document}
